\documentclass[letterpaper,11pt]{article}
\usepackage[margin=.67in,headheight=15pt,headsep=15pt,footskip=22pt]{geometry}
\usepackage[T1]{fontenc}\usepackage{lmodern,amsmath,amssymb,array,fancyhdr,hyperref}
\hypersetup{colorlinks=true,urlcolor=blue}
\renewcommand{\familydefault}{\sfdefault}\setlength{\parindent}{0pt}\setlength{\parskip}{6pt}
\pagestyle{fancy}\fancyhf{}\fancyhead[L]{Week 94 / Calibration lab / Facilitator guide}
\fancyfoot[L]{\scriptsize Bellingham Math Circle / Week 94 / W94-F-v2 / Unpiloted}\fancyfoot[R]{\small\thepage}
\pdfinfoomitdate=1\pdftrailerid{}\pdfsuppressptexinfo=15
\newcommand{\topic}[1]{\par\medskip\textbf{#1}\par}
\begin{document}
\topic{Mathematical overview: what can a measurement identify?}
\textbf{Shared offset.} In an exact addition machine, $y=x+c$. Thus $y_2-y_1=x_2-x_1$: the same offset cancels. Unknown lengths and an unknown offset can be changed together, $x\mapsto x+t$, $c\mapsto c-t$, without changing any reports, provided the changed values remain allowed. Repeating those reports cannot distinguish such worlds. A known reference $x$ identifies $c=y-x$. The middle packet restricts $c$ to $1,\ldots,6$; its particular pair of reports admits six worlds. The K--1 tower machine restricts $c$ to $1,2,3$, and towers of at least one cube coming out as 4 and 6 admit three worlds. Do not claim every possible finite-prior observation is ambiguous.

\textbf{Two-point calibration.} In the exact affine model $y=a+bx$, with $b>0$ and unrestricted real shift $a$, two distinct known inputs determine
\[
 b=\frac{y_2-y_1}{x_2-x_1},\qquad a=y_1-bx_1.
\]
One known input leaves infinitely many positive slopes: for any $b>0$, choose $a=y_1-bx_1$. The upper packet's earlier finite game has only 21 rules and is different: some single reports identify its rule, and a sufficiently long reference separates all 21.

\textbf{Bounded display error.} Let two exact references be separated by $D>0$ and let the displayed-report difference be $\Delta$. If each display error is at most $e$, the true difference lies in $[\Delta-2e,\Delta+2e]$. The possible slopes are
\[
 [(\Delta-2e)/D,(\Delta+2e)/D]\cap(0,\infty).
\]
Before positivity truncates it, the interval has width $4e/D$, with attainable endpoints. The largest possible width is therefore $4e/D$. Greater reference separation improves this \emph{worst-case} bound; particular observed positive-slope ranges may be shorter. If zero is excluded, a lower endpoint at zero is an infimum, not an attained smallest positive stretch. The printed numerical examples avoid that issue. Repeated readings are not assumed independent or centered; no confidence percentage or averaging-away guarantee follows.

\textbf{Testing is not assuming a model.} A third observation can refute an affine rule fitted to two references. Agreement at finitely many inputs cannot force an untested value when no functional form is promised. The distinction is between fitting within a stipulated family and establishing that the family describes the world.

\textbf{Discovery and readiness.} K--1 first measures intervals on a scale with no zero. Then it plays the same offset machine with cube towers: each tower comes out taller than its input by the extra, two outputs fit several pairs of inputs with one fixed difference, and one known tower picks out the pair. This needs counting to about 9 and towers stood side by side; no subtraction is required. Grades 2--3 uses small integer addition and subtraction to separate a fixed instrument error from unknown objects. Grades 4--5 first plays a finite rule game, then uses differences, ratios and fractions; Problems 6--7 are a readiness-dependent return visit. Physical examples suggest the results; complete candidate lists, equal-difference arguments and two-endpoint witnesses explain them. Formal algebra is adult background, not a prerequisite to start.
\newpage
\topic{Preparation and a manageable first route}
Eleven children at three fixed tables (two kindergartners and two first graders; four third graders; two fourth graders and a fifth grader), with one adult anchored to each table. Prepare five pair kits; the third upper child can be keeper or referee and rotate. Per kit: one spare first student page, cut into the four strips A,B,C,D; keep its twelve-square row uncut. Also supply 12 unit cubes or counters, an opaque cover (a standing folder), pencil and blank paper/whiteboard. The two K--1 kits replace the 12 counters with about 30 interlocking cubes (any brand or size) and add a crayon: Problems 4--6 build and colour towers. Add a ruler and scissors for the adult. The paper unit is \textbf{12 mm}: lengths A,B,C,D are 24,36,48,72 mm, each 12 mm wide. The K--1 rulers and paper row share that scale. Cubes are counting units in the machine game; no commercial cube-to-print fit is claimed. Build middle hidden strips one at a time. Longer known references can be measured on the paper row or drawn on spare paper.

Print single-sided US Letter at Actual Size. Use spare copies for cutting so each child's work page remains intact. Twenty strip rectangles and five unit rows serve the five kits; allow roughly 10--15 minutes for printing, cutting and checking. Measure one printed unit before use. These preparation estimates, physical fit and procedures have \textbf{not been rehearsed}.

\textbf{Launch together, about four minutes after movement.} Let everyone handle a strip and the unit row. Cut one square from a spare unit row and lay it beside ticks 8 and 9 of a printed ruler, as in the K--1 page 1 picture (or draw it on the board); show its two edge marks and its one-unit length. Then show the middle practice machine: one unit enters, one extra unit is added, report 2 leaves. The input length is known in this demonstration. Hide only the extra choice for a practice round; after a guess, reveal and check the whole operation. Keep the chosen rule unchanged within a round. Do not give a subtraction algorithm or the two-reference theorem at launch.

\textbf{About 35--40 minutes at tables.} K--1: handle and measure, then Problems 1--2 (about 15 minutes). Before Problem 4, play one round of the tower machine for the table with the page 3 picture: two cubes in, one secret extra cube added behind the folder, three cubes out. Problems 4--6 take about 20 minutes; Problem 3 suits a child who finishes early. Middle: play Problem 1, explore Problems 2--3, then the known-strip rows of Problem 4 and, if ready, Problem 5's game; 6--7 on a return visit. Upper: play 1--2, then choose 3--4 or return for them. Keep 5 and the whole last page for interested children or another meeting. Reserve final minutes for showing two different explanations of the same report. Finishing pages is not the aim.

\newpage
\topic{K--1 solutions, hints and extensions}
\textbf{1.} A=2, B=3, C=4, D=6 units wherever placed. The endpoint labels change together; the number of intervals does not. Hint: slide the uncut unit row alongside the strip, touching one space at a time. Do not count the tick marks as the length.

\textbf{2.} Both ends must remain at printed marks, as the shared rule states. A: 5--7 or 7--9; B: 4--7 or 7--10; C: 3--7 or 7--11; D: 7--13 only. D from 1--7 has the right length but is outside this printed board. These exhaust the choices because a straight strip has only two directions from its fixed end. Hint: turn the strip around its end at 7. Extension: choose another endpoint and investigate which strips fit on both sides.

\textbf{3.} A+C=D, with A and C in either order. There is one choice of two supplied pieces; B+B is unavailable because the kit has one B. Other distinct-piece sums are 5,8,7,9,10, none the length of an unused strip. Let physical covers establish the comparison first. Extension: make a new strip that can cover a different pair, and move the whole arrangement without breaking the match.

\textbf{Tower machine materials.} Each K--1 pair needs about 30 interlocking cubes, a standing folder as the cover, and a crayon. The keeper builds each tower that comes out in one colour, or mixed colours, but never the extra in a colour of its own: the finder should find the extra by comparing towers. Cube counts: Problem 5 with all three answers built at once uses 18 cubes, 28 with the two towers that came out.

\textbf{4. Play the tower machine (about 8 minutes).} The keeper chooses an extra of 1, 2 or 3 and keeps it for the round. For each tower the partner hands over, the keeper builds behind the folder a tower with the same number of cubes plus the extra, then shows it. One known tower is enough: stood beside a twin of the tower that went in, the tower that came out sticks up by the extra. Children often ask for several towers at first; afterwards ask ``How few towers did you need?'' so that one being enough becomes a finding. Hint: if the finder handed over their only tower, suggest building a twin to keep. Extension: send in two towers of different heights; how much taller is each output than its input, and how much taller is one output than the other?

\textbf{5. Every pair that could have gone in (about 7 minutes).} Extra 1: P 3, Q 5. Extra 2: P 2, Q 4. Extra 3: P 1, Q 3. The fourth slot stays empty. These are all, because each allowed extra has been used once; extra 0 is not allowed, so P and Q are never 4 and 6 themselves. P cannot be taller: Q is always 2 cubes taller, because taking the same extra off both towers leaves the 2 cubes Q stood above P. Hint: choose an extra and take that many cubes off the top of each tower that came out. Extension: if extras up to 5 were allowed, which new pairs would appear? None, since extra 4 would leave P with no cubes.

\textbf{6. A known tower picks the pair (about 6 minutes).} First machine: 1 cube in, 2 out, extra 1, so P 3 and Q 5. Second: 2 in, 5 out, extra 3, so P 1 and Q 3. Third: 3 in, 5 out, extra 2, so P 2 and Q 4. Each answer is one of Problem 5's pairs: the known tower chooses which. The second and third machines turn out the same 5-cube tower from different inputs, so the tower that went in matters as much as the one that came out. Hint: stand the known tower's input and output side by side. Extension: make a machine card for a partner, with a known tower in and out and P and Q coming out as 4 and 6.
\newpage
\topic{Grades 2--3 solutions and facilitation}
\textbf{1. Play the addition machine.} One exact known input and its report determine the extra. For example, length 3 reporting 8 means extra 5. A child may find the extra by matching units; subtraction notation can follow. The keeper chooses one extra for a whole round and reveals it afterwards. Hint: place the input units next to a row as long as the report and inspect what remains. Extension: compare two requested strips under the same extra; what changes by the same amount?

\textbf{2. Find the hidden lengths.} These are P and Q, \emph{not} the known cutouts A and B. The six allowed extras give:
\begin{center}\begin{tabular}{c|rrrrrr}
Extra&1&2&3&4&5&6\\\hline
P&6&5&4&3&2&1\\
Q&9&8&7&6&5&4
\end{tabular}\end{center}
Every pair fits the same reports 7 and 10. There are no others because every allowed extra has been used once. Allow pooled cases or an adult recorder. Hint: select one extra, then build each hidden length separately; do not require holding two cube rows at once. Extension: make another report pair with several explanations.

\textbf{3. What remains certain?} P cannot be longer; Q is always three units longer. Both reports include the same extra, so removing equal pieces preserves their three-unit difference. This is an explanation with objects, not merely noticing the six table rows. Extra 1 versus extra 6 gives a clear pair of different worlds. Extension: ask which facts change if each strip had its own unrelated extra.

\textbf{4. A known standard.} The known lengths 4,6,2 reporting 8 give extras 4,2,6 respectively. Hence hidden lengths reporting 9 and 12 are $(5,8)$, $(7,10)$, $(3,6)$. Each row is a separate machine; within a row the offset stays fixed. Build each hidden strip in turn. Hint: check the known strip against the report before interpreting either hidden strip. Extension: replace the standard while keeping the same machine; consistent standards agree on the extra.

\textbf{5. One requested reference.} The keeper chooses extra $c\in\{1,\ldots,6\}$ and writes hidden lengths $7-c,10-c$ under the cover. A known requested strip of length $r$ reports $r+c$, revealing $c$. The partner builds the two inferred strips in turn, then both check against the covered lengths. The keeper must not change $c$ after hearing the requested standard. Hint: first rehearse an open round. Extension: ask why repeating the two unknown-strip reports alone gives no new distinction.

\textbf{6. Detect and repair.} Box I has extras 3,3,4: change the last report 10 to 9. II has extra 6 throughout and is possible. III has extras 1,2,1: change the middle report 6 to 5. Each repair is unique because the other two correct pairs fix the extra. Hint: show one consistent box physically before discussing a repair. Extension: construct a box where one change cannot repair it.

\textbf{7. Invent a repair puzzle.} Choose one allowed extra, make three correct length-and-report pairs, then alter exactly one report. Give both inputs and reports. Two correct pairs identify the extra, so the wrong report is recoverable. Check the original rule on reveal. A changed input would be a different task. Extension: two bad reports can leave the intended extra ambiguous; build a counterexample instead of asserting that a partner's failure proves impossibility.
\newpage
\topic{Grades 4--5: finite choices, calibration and model checks}
\textbf{1. Finite machine game.} There are $3\times7=21$ choices $(b,a)$ with copies $b=1,2,3$ and extra $a=0,\ldots,6$. Let children choose requests and discover how outputs separate possibilities. Two distinct exact known inputs always determine a rule. A lucky single report can suffice: input 6, report 24 forces $b=3,a=6$. If arbitrary known lengths can be built, input 7 separates all rules: outputs are 7--13, 14--20 or 21--27 according to copies. With \emph{only} the four supplied lengths 2,3,4,6 requested one strip at a time, no single request guarantees identification; each has some ambiguous reports. With the kit, B and C laid end to end make 7, and that one request identifies every rule (A and D make 8; any total of 7 or more works). State which materials a child chose rather than enforcing an unstated query restriction. Hint: represent each copy count by its range of possible reports.

\textbf{2. One reference, then another.} Input 2 reporting 6 leaves $(b,a)=(1,4),(2,2),(3,0)$. A second input $r\ne2$ gives $r+4,2r+2,3r$, all different: equality between any pair would require $r=2$. Input 3 gives 7,8,9. The final question has answer \emph{yes} in this finite game; use Problem 1's examples. Extension: find a shortest integer reference whose one report separates all 21 rules (7); the intervals for copies 1 and 2 overlap for every integer input at most 6.

\textbf{3. Positive real stretch, signed shift.} Stop and run the printed non-task fractional example: input 3 becomes $4\tfrac12$ after stretching, then $3\tfrac12$ after subtracting 1. This introduces a new family of machines. A has rule $y=2x+4$, with reports 4,12,24 at lengths 0,4,10. B has rule $y=x+6$, with reports 6,10,16. For each, compare the change in reports when the input increases by 4. Any other affine rule through both points has that same stretch and then the same shift, so there is no other. Hint: compare the two reference changes before asking for the report at zero. Extension: choose fractional references; their separation, not their location, identifies stretch.

\textbf{4. Repair exactly one report.} The two input steps both have length 2, while the report steps are 4 and 5, so the three pairs do not fit one affine rule. Keeping each pair of reports in turn gives exactly three repairs:
\begin{center}\begin{tabular}{c|c|c}
Changed input's report & Corrected report & Repaired rule\\\hline
2&6&$y=\tfrac52x+1$\\
4&$11\tfrac12$&$y=\tfrac94x+\tfrac52$\\
6&15&$y=2x+3$
\end{tabular}\end{center}
All stretches are positive. There are no other one-report repairs because each choice of two unchanged points determines its line. Hint: keep two reports still while testing the third. Extension: find a four-report list that cannot be repaired by changing one.

\textbf{5. A model is no longer promised.} Example table A for lengths 1--8 is $5,7,9,11,13,15,17,19$; table B changes only the final 19 to 20. Both fit all three given reports. To define rules beyond the tables, take A as $2x+3$ everywhere and B equal to A for $x<8$ but A+1 for $x\ge8$. They agree on every test below 8 yet disagree at 8. No regularity or affine assumption is given, so those tests cannot force the untested report. Hint: identify which earlier promise was removed. Extension: discuss which new assumption would restore an answer; an assumed affine rule is enough, an unchanged but otherwise arbitrary rule is not.
\newpage
\topic{Grades 4--5: noisy displays and experimental design}
\textbf{6. True report versus displayed report.} Rehearse the printed example using one unchanged input: length 3, fixed rule $2x+1$, true report 7, display error $+\tfrac12$, display $7\tfrac12$. Only the last transition is noisy. The input standards stay exact and the machine stays fixed across both observations.

For pair A, the true reports at 2 and 4 lie in $[7,9]$ and $[11,13]$. Their difference is between 2 and 6, across input separation 2. Thus stretch is in $[1,3]$. Endpoints are reached by $y=x+7$ (true reports 9,11) and $y=3x+1$ (7,13). Pair B has true-report intervals $[7,9]$ and $[23,25]$ across separation 8. Its stretch range is $[7/4,9/4]$, attained by $y=(7/4)x+11/2$ (9,23) and $y=(9/4)x+5/2$ (7,25). Every intermediate stretch is feasible by interpolating the endpoint rules. The exact reference values must not acquire their own errors in this task.

\textbf{Held hint ladder.} First ask what the displayed 8 permits the true report to be. If needed, draw both permitted report bands on a whiteboard. Only later ask which pair of band endpoints makes the true-report change smallest or largest. Children can reason with slopes as ``report units gained per input unit''; naming a parameter is optional. A calculator or adult arithmetic can support fractional endpoints while children choose the extremes and check them.

\textbf{7. Choose before seeing the data.} The reference pair 1 and 10 has separation 9, giving worst-case slope-range width $4/9$. Any closer pair has larger worst-case width $4/D$. One proof: the two true reports can independently move through bands of width 2; their difference can move through width 4. Dividing by the input separation gives the slope width. Positive stretch can truncate a particular observed interval but cannot improve the largest possible width; for example, a true stretch of 3 with exact displays avoids truncation for every allowed separation $D\ge1$, since its interval $[3-2/D,\,3+2/D]$ starts at or above 1. (A true stretch of 2 does not: at $D=1$ its interval is $[0,4]$.) For any positive separation $D$, a true stretch above $2/D$ avoids truncation. Hence the endpoint references minimize the worst-case bound over lengths from 1 to 10.

Ten repeats of one reference cannot give such a bound without additional assumptions. Even ten identical, exact readings $y$ at length $x$ fit every positive stretch $b$ by taking shift $a=y-bx$. Different noisy readings may better constrain the true value at \emph{that input}; they still leave stretch free. Do not infer independence, a mean-zero error, a known shift, or a probabilistic confidence guarantee. Hint: let the shift compensate for a changed stretch at the one known input. Extension: change the display-error allowance to $e$ and investigate how the worst-case width changes; then separately ask what happens when the standards themselves are uncertain.

\topic{Connections, sources and use record}
Week 6 asks which exact questions distinguish finite hidden states. Here a shared unknown instrument parameter creates observationally identical explanations. Week 51's common-piece cancellation supplies a useful connection, but calibration chooses additional measurements to identify the machine. Atlas AP-05 studies repeated-sampling coverage; this packet's final page is a deterministic bounded-error question with no confidence percentage.

The independently derived toy models are grounded in NIST/SEMATECH e-Handbook sections \href{https://www.itl.nist.gov/div898/handbook/mpc/section3/mpc361.htm}{2.3.6.1, Models for instrument calibration}, \href{https://www.itl.nist.gov/div898/handbook/mpc/section3/mpc363.htm}{2.3.6.3, Assumptions for instrument calibration}, and \href{https://www.itl.nist.gov/div898/handbook/mpc/section4/mpc45.htm}{2.4.5, Analysis of bias}. They do not reproduce NIST's regression procedure. The coordinator's teaching-source review of \emph{Math Circle by the Bay}, printed pp. viii--x, supports deep themes, manipulatives and flexible pacing; these particular activities are our new design.

All packets are unpiloted. Record which tasks children actually used, whether they kept the same rule within a round, whether they confused input with display, and what they chose to continue. Digital mathematics and rendering checks do not establish physical fit or classroom success.
\end{document}
